\input eplain
\pdfpagewidth=148mm
\pdfpageheight=210mm
\hsize=128mm
\vsize=190mm

\pdfhorigin=0pt
\pdfvorigin=0pt
\hoffset=10mm
\voffset=10mm
\noindent news \leaders\hbox to 10pt {\hss.\hss}  \hfill \xref{node1}
\smallskip
\noindent toc \leaders\hbox to 10pt {\hss.\hss}  \hfill \xref{node3}
\smallskip
\noindent 01 mathematic statements \leaders\hbox to 10pt {\hss.\hss}  \hfill \xref{node0}
\smallskip
\noindent 02 implications \leaders\hbox to 10pt {\hss.\hss}  \hfill \xref{node2}
\smallskip
\noindent 03 rules of logic \leaders\hbox to 10pt {\hss.\hss}  \hfill \xref{node4}
\smallskip
\noindent 04 proofs \leaders\hbox to 10pt {\hss.\hss}  \hfill \xref{node5}
\smallskip
\noindent 05 proofs about discrete structures \leaders\hbox to 10pt {\hss.\hss}  \hfill \xref{node6}
\smallskip
\noindent 01 problems and definitions \leaders\hbox to 10pt {\hss.\hss}  \hfill \xref{node9}
\smallskip
\noindent 02 trees \leaders\hbox to 10pt {\hss.\hss}  \hfill \xref{node10}
\smallskip
\noindent 03 planar graphs \leaders\hbox to 10pt {\hss.\hss}  \hfill \xref{node11}
\smallskip
\noindent 01 pascals arithmetical triangle \leaders\hbox to 10pt {\hss.\hss}  \hfill \xref{node12}
\smallskip
\noindent 02 combining outcomes \leaders\hbox to 10pt {\hss.\hss}  \hfill \xref{node13}
\smallskip
\noindent 01 sets \leaders\hbox to 10pt {\hss.\hss}  \hfill \xref{node8}
\smallskip
\noindent 02 functions \leaders\hbox to 10pt {\hss.\hss}  \hfill \xref{node7}
\smallskip
\vfill \break
{\medskip\noindent\bf\font\titlefont=cmss17 at 20pt\titlefont Timeline}\medskip
{\noindent \bf 2025-10-19:} 01 mathematic statements\smallskip
{\noindent \bf 2025-10-20:} news, 01 mathematic statements\smallskip
{\noindent \bf 2025-10-21:} 02 implications\smallskip
{\noindent \bf 2025-10-28:} news, toc\smallskip
{\noindent \bf 2025-10-30:} news\smallskip
{\noindent \bf 2025-11-04:} 03 rules of logic\smallskip
{\noindent \bf 2025-11-05:} 03 rules of logic\smallskip
{\noindent \bf 2025-11-06:} 04 proofs\smallskip
{\noindent \bf 2025-11-09:} news, 04 proofs, 05 proofs about discrete structures\smallskip
{\noindent \bf 2025-11-10:} 05 proofs about discrete structures\smallskip
{\noindent \bf 2025-11-11:} 05 proofs about discrete structures\smallskip
{\noindent \bf 2025-11-13:} news, 02 functions\smallskip
{\noindent \bf 2025-11-14:} 05 proofs about discrete structures, 02 functions\smallskip
{\noindent \bf 2025-11-15:} 05 proofs about discrete structures\smallskip
{\noindent \bf 2025-11-18:} 01 sets\smallskip
{\noindent \bf 2025-11-19:} news, 01 sets\smallskip
{\noindent \bf 2025-11-28:} news\smallskip
{\noindent \bf 2025-12-01:} 01 problems and definitions\smallskip
{\noindent \bf 2025-12-02:} news, 01 problems and definitions, 02 trees\smallskip
{\noindent \bf 2025-12-03:} news\smallskip
{\noindent \bf 2026-01-06:} 02 trees\smallskip
{\noindent \bf 2026-01-09:} 03 planar graphs\smallskip
{\noindent \bf 2026-01-10:} 03 planar graphs\smallskip
{\noindent \bf 2026-01-11:} news\smallskip
{\noindent \bf 2026-01-12:} 03 planar graphs\smallskip
{\noindent \bf 2026-01-13:} 03 planar graphs\smallskip
{\noindent \bf 2026-01-16:} 01 pascals arithmetical triangle\smallskip
{\noindent \bf 2026-01-17:} 01 pascals arithmetical triangle\smallskip
{\noindent \bf 2026-01-18:} 01 pascals arithmetical triangle\smallskip
{\noindent \bf 2026-01-19:} 02 combining outcomes\smallskip
{\noindent \bf 2026-01-21:} 02 combining outcomes\smallskip
{\noindent \bf 2026-01-22:} news, 02 combining outcomes\smallskip
{\noindent \bf 2026-01-26:} news, 02 combining outcomes\smallskip
{\noindent \bf 2026-01-27:} 02 combining outcomes\smallskip
\vfill \break
\xrdef{node0}
{\medskip\noindent\bf\font\titlefont=cmss17 at 20pt\titlefont 01 mathematic statements}\medskip
{\medskip\noindent\bf\font\titlefont=cmitt12 at 17pt\titlefont 2025-10-19}\medskip
{\smallskip\noindent\bf\font\titlefont=cmr12 at 13pt\titlefont $\bullet$ 21:06: 1.1 Mathematical Statements}\smallskip
{\medskip\noindent\bf\font\titlefont=cmitt12 at 17pt\titlefont 2025-10-20}\medskip
{\smallskip\noindent\bf\font\titlefont=cmr12 at 13pt\titlefont $\bullet$ 21:28: DMOI 1.1.4}\smallskip
\xrdef{node1}
{\medskip\noindent\bf\font\titlefont=cmss17 at 20pt\titlefont news}\medskip
{\medskip\noindent\bf\font\titlefont=cmitt12 at 17pt\titlefont 2025-10-20}\medskip
{\smallskip\noindent\bf\font\titlefont=cmr12 at 13pt\titlefont $\bullet$ 09:28: started reading this textbook yesterday}\smallskip
{\medskip\noindent\bf\font\titlefont=cmitt12 at 17pt\titlefont 2025-10-28}\medskip
{\smallskip\noindent\bf\font\titlefont=cmr12 at 13pt\titlefont $\bullet$ 17:55: Set up toc, make it suitable for tasks}\smallskip
{\smallskip\noindent\bf\font\titlefont=cmr12 at 13pt\titlefont $\bullet$ 19:10: Back to toc and planning}\smallskip
{\medskip\noindent\bf\font\titlefont=cmitt12 at 17pt\titlefont 2025-10-30}\medskip
{\smallskip\noindent\bf\font\titlefont=cmr12 at 13pt\titlefont $\bullet$ 10:38: Generating initial reading schedule for DMOI}\smallskip
Hoping to combine this with the sonilo schedule and build a unified agenda
\par

{\medskip\noindent\bf\font\titlefont=cmitt12 at 17pt\titlefont 2025-11-09}\medskip
{\smallskip\noindent\bf\font\titlefont=cmr12 at 13pt\titlefont $\bullet$ 15:35: New books! COAD, DMOI}\smallskip
{\medskip\noindent\bf\font\titlefont=cmitt12 at 17pt\titlefont 2025-11-13}\medskip
{\smallskip\noindent\bf\font\titlefont=cmr12 at 13pt\titlefont $\bullet$ 10:52: experiment: transcribing a page from DMOI that used my new graph system}\smallskip
This is a new handwritten knowledge graph system that divides the page into to parts: a list of nodes with a symbol, and a graph, that connects nodes (symbols only) with wires.
\par

I'm going to open up ed, and attemp to transcribe this. Making this up as I go. Timeboxing to 15 minutes.
\par

Interesting experiment. I got halfway done with a page in 15 minutes. Translating the text was tricky. The graph traversal I did left to write top to bottom. I used a ruler and another paper to limit visibility to only the nodes I needed to input.
\par

30-45 minutes processing time per page I'm guessing? Things that have more math symbols will take more time.
\par

parsing this data will come next. it needs to be a very permissive parser.
\par

{\medskip\noindent\bf\font\titlefont=cmitt12 at 17pt\titlefont 2025-11-19}\medskip
{\smallskip\noindent\bf\font\titlefont=cmr12 at 13pt\titlefont $\bullet$ 19:52: attempting to transcribe another page of notes}\smallskip
It's a non-trivial effort still. but, I have a feeling I'll get faster at it. The parser does seem to work!
\par

I've added my parsing utility to "tools", and my transcription of 05H to a new folder.
\par

{\medskip\noindent\bf\font\titlefont=cmitt12 at 17pt\titlefont 2025-11-28}\medskip
{\smallskip\noindent\bf\font\titlefont=cmr12 at 13pt\titlefont $\bullet$ 09:17: addressing the agenda overflow}\smallskip
As of today, there are 26 tasks total.
\par

9 of those tasks are related to sonilo. I am stuck. Memory allocators have taken time. Ilo has taken time. Trying to get the moving parts right is taking time. Sonilo ilo interpreter is about 26 days overdue, that's the most overdue task.
\par

COAD also has 9 tasks, surprisingly. I guess this makes a little more sense. I'll have gone about 10 days without doing any reading. The most overdue task is 13 days, so I'm about 3-5 days overdue factoring in this trip.
\par

DMOI is 11 days overdue for the latest. That means I'm about 1-3 days overdue factoring vacation. Not too bad.
\par

I am going to add in a way to get the max day for selected groups, then shift sonilo and COAD back 10 days. Once I have this "max days" functionality, I should be able to calculate a "due date" from the agenda.
\par

{\medskip\noindent\bf\font\titlefont=cmitt12 at 17pt\titlefont 2025-12-02}\medskip
{\smallskip\noindent\bf\font\titlefont=cmr12 at 13pt\titlefont $\bullet$ 15:17: spreading out reading stuff more, agenda is too piled up}\smallskip
{\medskip\noindent\bf\font\titlefont=cmitt12 at 17pt\titlefont 2025-12-03}\medskip
{\smallskip\noindent\bf\font\titlefont=cmr12 at 13pt\titlefont $\bullet$ 11:02: shifting goalposts again for sonilo project, dmoi}\smallskip
I need these tasks to go down, so I'm redistributing the days. I'm disappointed at velocity, but at least I'm measuring it. the number is now 180 days
\par

stretching the reading for dmoi so only 2 tasks appear right now the reading takes a while. I am a thorough (read: slow) reader.
\par

{\medskip\noindent\bf\font\titlefont=cmitt12 at 17pt\titlefont 2026-01-11}\medskip
{\smallskip\noindent\bf\font\titlefont=cmr12 at 13pt\titlefont $\bullet$ 14:56: some DMOI transcription}\smallskip
a new transcription group for DMOI was created!
\par

nodes, edges, names created. figure is a WIP. still hadn't fully gotten the hang of drawing these things in a way that was easy to transcribe.
\par

{\smallskip\noindent\bf\font\titlefont=cmr12 at 13pt\titlefont $\bullet$ 15:34: boilerplate transcriptions to get ink 05v to work}\smallskip
{\medskip\noindent\bf\font\titlefont=cmitt12 at 17pt\titlefont 2026-01-22}\medskip
{\smallskip\noindent\bf\font\titlefont=cmr12 at 13pt\titlefont $\bullet$ 10:02: increasing DMOI days}\smallskip
This is just going to take a while, apparently
\par

Now set to 250 days.
\par

{\medskip\noindent\bf\font\titlefont=cmitt12 at 17pt\titlefont 2026-01-26}\medskip
{\smallskip\noindent\bf\font\titlefont=cmr12 at 13pt\titlefont $\bullet$ 10:59: delayed a month, spread increased to 200 days}\smallskip
I'll probably need to revisit this anyways. School might slow this down.
\par

\xrdef{node2}
{\medskip\noindent\bf\font\titlefont=cmss17 at 20pt\titlefont 02 implications}\medskip
{\medskip\noindent\bf\font\titlefont=cmitt12 at 17pt\titlefont 2025-10-21}\medskip
{\smallskip\noindent\bf\font\titlefont=cmr12 at 13pt\titlefont $\bullet$ 13:20: 1.2 Implications}\smallskip
\xrdef{node3}
{\medskip\noindent\bf\font\titlefont=cmss17 at 20pt\titlefont toc}\medskip
{\medskip\noindent\bf\font\titlefont=cmitt12 at 17pt\titlefont 2025-10-28}\medskip
{\smallskip\noindent\bf\font\titlefont=cmr12 at 13pt\titlefont $\bullet$ 17:55: Set up toc, make it suitable for tasks}\smallskip
{\smallskip\noindent\bf\font\titlefont=cmr12 at 13pt\titlefont $\bullet$ 19:10: Back to toc and planning}\smallskip
\xrdef{node4}
{\medskip\noindent\bf\font\titlefont=cmss17 at 20pt\titlefont 03 rules of logic}\medskip
{\medskip\noindent\bf\font\titlefont=cmitt12 at 17pt\titlefont 2025-11-04}\medskip
{\smallskip\noindent\bf\font\titlefont=cmr12 at 13pt\titlefont $\bullet$ 17:03: 1.3 Rules of logic}\smallskip
{\medskip\noindent\bf\font\titlefont=cmitt12 at 17pt\titlefont 2025-11-05}\medskip
{\smallskip\noindent\bf\font\titlefont=cmr12 at 13pt\titlefont $\bullet$ 14:29: DMOI 1.3 rules of logic reading}\smallskip
{\smallskip\noindent\bf\font\titlefont=cmr12 at 13pt\titlefont $\bullet$ 16:30: DMOI 1.3 continued}\smallskip
\xrdef{node5}
{\medskip\noindent\bf\font\titlefont=cmss17 at 20pt\titlefont 04 proofs}\medskip
{\medskip\noindent\bf\font\titlefont=cmitt12 at 17pt\titlefont 2025-11-06}\medskip
{\smallskip\noindent\bf\font\titlefont=cmr12 at 13pt\titlefont $\bullet$ 12:50: 1.4 proofs}\smallskip
{\medskip\noindent\bf\font\titlefont=cmitt12 at 17pt\titlefont 2025-11-09}\medskip
{\smallskip\noindent\bf\font\titlefont=cmr12 at 13pt\titlefont $\bullet$ 09:08: DMOI 1.4}\smallskip
\xrdef{node6}
{\medskip\noindent\bf\font\titlefont=cmss17 at 20pt\titlefont 05 proofs about discrete structures}\medskip
{\medskip\noindent\bf\font\titlefont=cmitt12 at 17pt\titlefont 2025-11-09}\medskip
{\smallskip\noindent\bf\font\titlefont=cmr12 at 13pt\titlefont $\bullet$ 14:15: 1.5 proofs about discrete structures}\smallskip
{\medskip\noindent\bf\font\titlefont=cmitt12 at 17pt\titlefont 2025-11-10}\medskip
{\smallskip\noindent\bf\font\titlefont=cmr12 at 13pt\titlefont $\bullet$ 12:17: DMOI 1.5}\smallskip
{\smallskip\noindent\bf\font\titlefont=cmr12 at 13pt\titlefont $\bullet$ 14:20: DMOI 1.5 continued}\smallskip
{\medskip\noindent\bf\font\titlefont=cmitt12 at 17pt\titlefont 2025-11-11}\medskip
{\smallskip\noindent\bf\font\titlefont=cmr12 at 13pt\titlefont $\bullet$ 21:12: DMOI 1.5}\smallskip
{\medskip\noindent\bf\font\titlefont=cmitt12 at 17pt\titlefont 2025-11-14}\medskip
{\smallskip\noindent\bf\font\titlefont=cmr12 at 13pt\titlefont $\bullet$ 15:33: DMOI 1.5: almost done reading the chapter}\smallskip
{\medskip\noindent\bf\font\titlefont=cmitt12 at 17pt\titlefont 2025-11-15}\medskip
{\smallskip\noindent\bf\font\titlefont=cmr12 at 13pt\titlefont $\bullet$ 12:55: DMOI 1.5 wrapping up}\smallskip
\xrdef{node7}
{\medskip\noindent\bf\font\titlefont=cmss17 at 20pt\titlefont 02 functions}\medskip
{\medskip\noindent\bf\font\titlefont=cmitt12 at 17pt\titlefont 2025-11-13}\medskip
{\smallskip\noindent\bf\font\titlefont=cmr12 at 13pt\titlefont $\bullet$ 14:28: DMOI 5.2 functions}\smallskip
{\medskip\noindent\bf\font\titlefont=cmitt12 at 17pt\titlefont 2025-11-14}\medskip
{\smallskip\noindent\bf\font\titlefont=cmr12 at 13pt\titlefont $\bullet$ 13:40: what a waste of time}\smallskip
I read this chapter in preparation for the quiz the following day. In hindsight, it did not help with the quiz.
\par

Between re-reading the zybooks chapter that the quiz was supposed to be on, and reading the DMOI chapter that was related to the material in the zybooks chapter, I logged 3+ hours. over the course of 2 days. The fact that I showed up and none of it helped makes me really upset. The quiz had nothing from the chapter it was all off book. It was open book too, and I still probably failed.
\par

I don't think this was my fault. The quiz literally had nothing from the textbook, and used completely different vocabulary. Considering that up to this point, this has been a rote math class, I was pretty shocked by the quiz.
\par

I want my hours of studying back. What a waste.
\par

\xrdef{node8}
{\medskip\noindent\bf\font\titlefont=cmss17 at 20pt\titlefont 01 sets}\medskip
{\medskip\noindent\bf\font\titlefont=cmitt12 at 17pt\titlefont 2025-11-18}\medskip
{\smallskip\noindent\bf\font\titlefont=cmr12 at 13pt\titlefont $\bullet$ 16:10: DMOI 5.1: Sets}\smallskip
{\smallskip\noindent\bf\font\titlefont=cmr12 at 13pt\titlefont $\bullet$ 20:25: DMOI 5.1 continued}\smallskip
{\medskip\noindent\bf\font\titlefont=cmitt12 at 17pt\titlefont 2025-11-19}\medskip
{\smallskip\noindent\bf\font\titlefont=cmr12 at 13pt\titlefont $\bullet$ 19:52: attempting to transcribe another page of notes}\smallskip
It's a non-trivial effort still. but, I have a feeling I'll get faster at it. The parser does seem to work!
\par

I've added my parsing utility to "tools", and my transcription of 05H to a new folder.
\par

\xrdef{node9}
{\medskip\noindent\bf\font\titlefont=cmss17 at 20pt\titlefont 01 problems and definitions}\medskip
{\medskip\noindent\bf\font\titlefont=cmitt12 at 17pt\titlefont 2025-12-01}\medskip
{\smallskip\noindent\bf\font\titlefont=cmr12 at 13pt\titlefont $\bullet$ 15:06: started chapter 2, 2.1}\smallskip
{\smallskip\noindent\bf\font\titlefont=cmr12 at 13pt\titlefont $\bullet$ 21:10: More 2.1}\smallskip
{\medskip\noindent\bf\font\titlefont=cmitt12 at 17pt\titlefont 2025-12-02}\medskip
{\smallskip\noindent\bf\font\titlefont=cmr12 at 13pt\titlefont $\bullet$ 14:00: 2.1 wrap-up}\smallskip
\xrdef{node10}
{\medskip\noindent\bf\font\titlefont=cmss17 at 20pt\titlefont 02 trees}\medskip
{\medskip\noindent\bf\font\titlefont=cmitt12 at 17pt\titlefont 2025-12-02}\medskip
{\smallskip\noindent\bf\font\titlefont=cmr12 at 13pt\titlefont $\bullet$ 21:30: 2.2 Trees}\smallskip
{\medskip\noindent\bf\font\titlefont=cmitt12 at 17pt\titlefont 2026-01-06}\medskip
{\smallskip\noindent\bf\font\titlefont=cmr12 at 13pt\titlefont $\bullet$ 14:00: DMOI 2.2.3 spanning trees}\smallskip
spanning trees, minimum spanning trees, parent, child, grandparent, grandchild, ancestor, descendant, rooted trees.
\par

done with this section. I'm 39 days overdue as of today.
\par

\xrdef{node11}
{\medskip\noindent\bf\font\titlefont=cmss17 at 20pt\titlefont 03 planar graphs}\medskip
{\medskip\noindent\bf\font\titlefont=cmitt12 at 17pt\titlefont 2026-01-09}\medskip
{\smallskip\noindent\bf\font\titlefont=cmr12 at 13pt\titlefont $\bullet$ 16:17: 2.3.2 planar graphs}\smallskip
planar graphs are graphs where it's possible to draw without crossing edges. drawn verts and edges divdide the graph into regions called faces. eulers formula for planar graphs, deals with the relationship between edges, verts, and faces, which is a constant. a proof accompanies this.
\par

{\medskip\noindent\bf\font\titlefont=cmitt12 at 17pt\titlefont 2026-01-10}\medskip
{\smallskip\noindent\bf\font\titlefont=cmr12 at 13pt\titlefont $\bullet$ 15:52: 2.3.3 non-planar graphs}\smallskip
non-planar graphs means there are too many edges and too few verts for it to be a planar graph, and some edges need to intersect. The smallest graph for this is is $K_5$, with a proof by contradiction using Euler's Formula for Planar graphs. Another graph examined is $K_{3, 3}$, with a similar proof.
\par

{\medskip\noindent\bf\font\titlefont=cmitt12 at 17pt\titlefont 2026-01-12}\medskip
{\smallskip\noindent\bf\font\titlefont=cmr12 at 13pt\titlefont $\bullet$ 14:34: DMOI 2.3.4 polyhedra}\smallskip
polyhedron shapes, like planar graphs, are made up of vertices, edges, and faces. convex polyhedron can be represented as planar graphs. proof that there are only 5 regular polyhedra (have not studied this one yet)
\par

{\medskip\noindent\bf\font\titlefont=cmitt12 at 17pt\titlefont 2026-01-13}\medskip
{\smallskip\noindent\bf\font\titlefont=cmr12 at 13pt\titlefont $\bullet$ 11:55: DMOI: polyhydra proof initial outline}\smallskip
Just enough time to outline the 4 cases.
\par

{\smallskip\noindent\bf\font\titlefont=cmr12 at 13pt\titlefont $\bullet$ 14:36: DMOI polyhdra proof}\smallskip
\xrdef{node12}
{\medskip\noindent\bf\font\titlefont=cmss17 at 20pt\titlefont 01 pascals arithmetical triangle}\medskip
{\medskip\noindent\bf\font\titlefont=cmitt12 at 17pt\titlefont 2026-01-16}\medskip
{\smallskip\noindent\bf\font\titlefont=cmr12 at 13pt\titlefont $\bullet$ 17:18: 3.1 pascale's arithmetical triangle}\smallskip
chapter beings talking about chess moves, and somehow this is related to pizza toppings? Blaise Pascale, triangular array of numbers (figurate numbers). Pascal's triangle can be defined as every number sum of 2 above it. lattice paths, bit strings, subsets. Left off at notation.
\par

{\medskip\noindent\bf\font\titlefont=cmitt12 at 17pt\titlefont 2026-01-17}\medskip
{\smallskip\noindent\bf\font\titlefont=cmr12 at 13pt\titlefont $\bullet$ 14:11: wrap up 3.1.1, work on 3.1.2}\smallskip
Notation: "n choose k" are coordinates of pascal's triangle, row n column k. integer lattice, lattice path. using pascal's triangle to find number of lattice paths between two points. This was actually the dynamic programming assignment I did for cs342 "minCostPath".
\par

{\medskip\noindent\bf\font\titlefont=cmitt12 at 17pt\titlefont 2026-01-18}\medskip
{\smallskip\noindent\bf\font\titlefont=cmr12 at 13pt\titlefont $\bullet$ 13:18: 3.1.3 bit strings, 3.1.4 subsets and pizzas, 3.1.5 algebra?}\smallskip
3.1.3: bit strings, weights, lengths, it relates to lattice paths.
\par

3.1.4: counting subsets, it relates to bitstrings weights/lengths. choosing pizza toppings. "n choose k" means choosing k things from a set of n total things.
\par

3.1.5: binomials, expansion. binomial theorem. binomial coefficients are pascal's triangle numbers
\par

\xrdef{node13}
{\medskip\noindent\bf\font\titlefont=cmss17 at 20pt\titlefont 02 combining outcomes}\medskip
{\medskip\noindent\bf\font\titlefont=cmitt12 at 17pt\titlefont 2026-01-19}\medskip
{\smallskip\noindent\bf\font\titlefont=cmr12 at 13pt\titlefont $\bullet$ 15:55: 3.2 combining outcomes (3.2.1)}\smallskip
section preview. combinatorics. dealing with how things combine. a deck of cards. pick 2 cards, how many combos with first red, second face?
\par

{\smallskip\noindent\bf\font\titlefont=cmr12 at 13pt\titlefont $\bullet$ 17:05: more 3.2 today (3.2.2)}\smallskip
how to ask about counting problems? asking for number of elements in set of outcomes. how to combine outcomes? combine the set of outcomes or combine the outcomes in the sets.
\par

{\medskip\noindent\bf\font\titlefont=cmitt12 at 17pt\titlefont 2026-01-21}\medskip
{\smallskip\noindent\bf\font\titlefont=cmr12 at 13pt\titlefont $\bullet$ 14:40: DMOI 3.2.3: sum/product principles}\smallskip
sum principle, "subtraction principle", multiple events, multiplication leads to product principle. outcomes in product principle as trees with outcomes as levels, leaves of trees outcome. using product principle to define sets of outcomes that first need to be counted carefully
\par

a fair bit of example problems to study here.
\par

{\smallskip\noindent\bf\font\titlefont=cmr12 at 13pt\titlefont $\bullet$ 15:15: More 3.2.3}\smallskip
{\medskip\noindent\bf\font\titlefont=cmitt12 at 17pt\titlefont 2026-01-22}\medskip
{\smallskip\noindent\bf\font\titlefont=cmr12 at 13pt\titlefont $\bullet$ 10:35: 3.2.3 continued}\smallskip
product principle generalized to more than 2 events. applying product principle to functions (aka "rigoursly defined mathematical objects").
\par

{\smallskip\noindent\bf\font\titlefont=cmr12 at 13pt\titlefont $\bullet$ 11:18: 3.2.3, 3.2.4}\smallskip
3.2.4: combining principles. case study on red/face deck of cards problems. how to set up the outcomes?
\par

{\medskip\noindent\bf\font\titlefont=cmitt12 at 17pt\titlefont 2026-01-26}\medskip
{\smallskip\noindent\bf\font\titlefont=cmr12 at 13pt\titlefont $\bullet$ 10:58: this might be the last chapter I read for a while}\smallskip
It looks like I'm pretty close to wrapping this one up.
\par

{\medskip\noindent\bf\font\titlefont=cmitt12 at 17pt\titlefont 2026-01-27}\medskip
{\smallskip\noindent\bf\font\titlefont=cmr12 at 13pt\titlefont $\bullet$ 10:48: DMOI 3.2 wrap-up}\smallskip
This finishes up how to think about 2-card problem using product principle and sum principle. I'm not going to try to summarize it more than that. The numbers were kind of given, and I haven't had to chance to really work out the numbers and fully trust them.
\par

\bye
